% Full proof for the exact statement in
% p09p2_reduction_statement_exact_20260825.tex.

\paragraph{Imported inputs.}
The proof uses the following statements with exactly the hypotheses recorded
here.
\begin{enumerate}
\item Gubler--K\"unnemann work over an algebraically closed, complete,
nontrivially valued non-Archimedean field.  Their Proposition~8.11
identifies formal and piecewise linear metrics; their Proposition~8.13
identifies algebraic and formal metrics on line bundles over proper
varieties; their Definition~9.9 and \S9.13 define ordinary delta-metrics in
algebraic tropical charts; their \S9.14 gives isometry invariance, tensor
and dual closure, and descent from a positive tensor power; and their
Remark~9.16 makes piecewise linear, formal, and algebraic metrics
delta-metrics.  Their Example~8.16 states that the normalized canonical
metric of an algebraically trivial line bundle has an algebraic positive
tensor power.

\item The algebraicity statement used in
\cite[Example~8.16]{GK17} is \cite[Example~3.7]{Gub10CM}.  It is used only
in the form quoted by Gubler--K\"unnemann.

\item Chambert--Loir--Ducros,
\cite[Lemmas~19.1.4 and 19.1.6]{CLD25}, give compatibility of cubical
metrics with tensor products, duals, and homomorphism pullbacks, existence
and uniqueness up to a positive scalar, uniqueness after normalizing a
rigidification, and the exponent-one characterization for an
antisymmetric line bundle.

\item Milne, \cite[Proposition~8.4 and Remark~8.5]{milneAV}, gives the
characterization of \(\operatorname{Pic}^0(A)\) used below and inversion as
duality on \(\operatorname{Pic}^0(A)\); his
\cite[Remark~10.19]{milneAV} states that inversion acts trivially on
\(\operatorname{NS}(A)\).

\item The Stacks Project, \cite[Lemma~39.9.2]{stacks-project}, gives an
ample line bundle on an abelian variety.  Its Lemma~28.27.5
(Tag~0890) preserves ampleness under the tensor product used below, and
its Lemma~29.38.7 (Tag~0892) gives
\(S\otimes H^{\otimes m}\) ample for all sufficiently large \(m\) when
\(H\) is ample.
\end{enumerate}

\paragraph{Metric recognition and closure rules.}
The proof uses the following meanings and no identifications between them
without a cited recognition result.  An algebraic metric on a line bundle
over a proper variety is induced by a flat proper finite-type
\(K^\circ\)-model and a line bundle on that model.  A formal metric is
induced by an admissible formal \(K^\circ\)-model.  On a proper variety,
algebraic and formal metrics coincide by
\cite[Proposition~8.13]{GK17}; formal and piecewise linear metrics coincide
by \cite[Proposition~8.11]{GK17}.  A piecewise linear metric is an ordinary
delta-metric by \cite[\S9.13 and Remark~9.16]{GK17}.

For comparison with the local definition, a piecewise smooth weight
\(\phi\) in an ordinary tropical chart is a delta-weight when its facewise
Hessian
\[
                 \mathrm d'_{\mathrm P}\mathrm d''_{\mathrm P}\phi
\]
is induced by a codimension-zero generalized delta-preform on the ambient
tropical vector space.  This is the condition in
\cite[Definition~9.9 and \S9.13]{GK17}.  In particular, a constant real
weight is smooth and is a delta-weight, whether or not the constant lies in
the value group of \(K\).

We shall repeatedly use the exact closure assertions of
\cite[\S9.14]{GK17}: ordinary delta-metrics are invariant under isometry,
are closed under tensor product and duality, and satisfy
\begin{equation}
\label{eq:p09p2-root-rule}
 \overline M^{\otimes N}\text{ is an ordinary delta-metric for some }N\geq1
 \quad\Longrightarrow\quad
 \overline M\text{ is an ordinary delta-metric}.
\end{equation}
Thus \eqref{eq:p09p2-root-rule} is a theorem about the already given metric
on \(M\); it neither asserts nor requires the existence of an \(N\)-th root
of an unrelated line bundle or metric.

\begin{proof}
\cref{itm:p09p2-pic0}
Choose a rigidification \(\varepsilon\in L(0)\setminus\{0\}\).  Multiplying
the given metric by a positive constant, if necessary, gives a cubical
metric \(\|\cdot\|_0\) for which \(\|\varepsilon\|_0=1\).  Constant
rescaling preserves the cubical condition: the total exponent with which a
constant enters the metric on \eqref{eq:p09p2-d3} is
\[
        \sum_{I\subseteq\{1,2,3\}}(-1)^{\# I+1}=0.
\]

Since \(L\in\operatorname{Pic}^0(A)\), one has
\(\iota^*L\simeq L^\vee\).  Indeed, the characterization of
\(\operatorname{Pic}^0(A)\) in \cite[Proposition~8.4]{milneAV} gives
\(m^*L\simeq p_1^*L\otimes p_2^*L\) on \(A\times A\); pulling this
isomorphism back by \(x\mapsto(-x,x)\) yields
\(\iota^*L\otimes L\simeq\mathcal O_A\).  The normalized cubical metric
\(\|\cdot\|_0\) is therefore the normalized canonical metric used for an
algebraically trivial line bundle in \cite[Example~8.16]{GK17}; equivalently,
it is characterized by the isometric relation
\([n]^*\overline L_0\simeq\overline L_0^{\otimes n}\) for \(n\geq2\), as in
\cite[Lemmas~19.1.4 and 19.1.6]{CLD25}.

By \cite[Example~8.16]{GK17}, whose algebraicity input is
\cite[Example~3.7]{Gub10CM}, there is an integer \(N\geq1\) for which
\(\overline L_0^{\otimes N}\) is an algebraic metric.  Because \(A\) is
proper, the recognition chain recorded above gives
\[
 \overline L_0^{\otimes N}\text{ algebraic}
 \Longrightarrow
 \overline L_0^{\otimes N}\text{ formal}
 \Longrightarrow
 \overline L_0^{\otimes N}\text{ piecewise linear}
 \Longrightarrow
 \overline L_0^{\otimes N}\text{ an ordinary delta-metric}.
\]
Applying \eqref{eq:p09p2-root-rule} proves that \(\overline L_0\) is an
ordinary delta-metric.

Return to the original cubical metric.  It has the form
\(\overline L=\overline L_0\otimes\overline{\mathcal O}_{c}\), where
\(\overline{\mathcal O}_{c}\) is the trivial line bundle with a constant
positive norm.  Its local weight is constant, hence smooth and an ordinary
delta-weight.  Tensor closure now proves \cref{itm:p09p2-pic0} for the
original, arbitrarily normalized cubical metric.

\cref{itm:p09p2-split}
The operations in \eqref{eq:p09p2-plus-minus} preserve cubical metrics.
This follows directly from \eqref{eq:p09p2-d3}: \(D_3\) commutes with tensor
products and duals, and for a homomorphism \(u:A\to A\) one has
\(D_3(u^*L)=(u^3)^*D_3(L)\), compatibly with the distinguished
trivializations.  This is also the tensor, dual, and pullback compatibility
of \cite[Lemmas~19.1.4 and 19.1.6]{CLD25}.

Because \(\iota^2=\operatorname{id}_A\), the definitions give
\[
 \iota^*L_+
   =\iota^*L\otimes L\simeq L_+,
 \qquad
 \iota^*L_-
   =\iota^*L\otimes L^\vee\simeq L_-^\vee.
\]
Moreover, \(\iota^*\) acts as the identity on
\(\operatorname{NS}(A)=\operatorname{Pic}(A)/\operatorname{Pic}^0(A)\),
by \cite[Remark~10.19]{milneAV}.  Hence
\[
 [L_-]_{\operatorname{NS}}
   =[L]_{\operatorname{NS}}-[\iota^*L]_{\operatorname{NS}}=0,
\]
so \(L_-\in\operatorname{Pic}^0(A)\).  Finally, evaluation for the pullback
metric gives the isometry
\[
 \begin{aligned}
 \overline L_+\otimes\overline L_-
  &=(\overline L\otimes\iota^*\overline L)
    \otimes(\overline L\otimes(\iota^*\overline L)^\vee)\\
  &\simeq \overline L^{\otimes2},
 \end{aligned}
\]
which is \eqref{eq:p09p2-square-split}.  The rigidifications on
\(L_+\), \(L_-\), and \(L^{\otimes2}\) are the induced ones, and there is no
normalization scalar in this isometry because all four metrics are the
directly induced tensor, pullback, and dual metrics from
\eqref{eq:p09p2-plus-minus}.  In particular, when
\(\|\varepsilon\|_L=1\), the induced rigidifications on \(L_+\), \(L_-\),
and \(L^{\otimes2}\) also have norm one and the displayed isometry is an
isometry in the normalized rigidified category.

By \cref{itm:p09p2-pic0}, \(\overline L_-\) is an ordinary delta-metric.  If
\(\overline L_+\) is an ordinary delta-metric, tensor closure and
\eqref{eq:p09p2-square-split} make \(\overline L^{\otimes2}\) an ordinary
delta-metric, and \eqref{eq:p09p2-root-rule} makes \(\overline L\) one.  In
the reverse direction, if \(\overline L\) is an ordinary delta-metric, then
\[
       \overline L_+
       \simeq
       \overline L^{\otimes2}\otimes\overline L_-^\vee
\]
is an ordinary delta-metric by tensor and dual closure.  This proves
\eqref{eq:p09p2-delta-equivalence}.

\cref{itm:p09p2-ample-reduction}
The implication
\cref{itm:p09p2-all}\(\Rightarrow\)\cref{itm:p09p2-ample-symmetric} is
immediate.  Assume \cref{itm:p09p2-ample-symmetric}.  We first prove that
every cubical metric \(\overline S\) on every symmetric line bundle \(S\)
is an ordinary delta-metric.

An abelian variety is projective by
\cite[Lemma~39.9.2]{stacks-project}, so choose an ample line bundle \(H_0\)
on \(A\), choose a rigidification, and give it its normalized cubical metric
\(\overline H_0\), whose existence and uniqueness are
\cite[Lemmas~19.1.4 and 19.1.6]{CLD25}.  Set
\[
                 \overline H
        :=\overline H_0\otimes\iota^*\overline H_0.
\]
The preceding compatibility shows that \(\overline H\) is cubical, and
\(\iota^*H\simeq H\).  Pullback by the automorphism \(\iota\) preserves
ampleness, and the tensor product of ample line bundles is ample by
\cite[Lemma~28.27.5, Tag~0890]{stacks-project}; hence \(H\) is ample.

For all sufficiently large integers \(m\), the line bundle
\begin{equation}
\label{eq:p09p2-ample-twist}
                         S_m:=S\otimes H^{\otimes m}
\end{equation}
is ample.  This is
\cite[Lemma~29.38.7, Tag~0892]{stacks-project}, applied to the identity
morphism of \(A\): every invertible sheaf is relatively ample for the
identity morphism, and \(H\) is ample on the target.  Both \(H\) and \(S_m\)
are symmetric.  Give \(S_m\) the induced rigidification and cubical metric
\(\overline S_m=\overline S\otimes\overline H^{\otimes m}\).  The assumed
statement \cref{itm:p09p2-ample-symmetric} applies to \(\overline H\) and
\(\overline S_m\), so both are ordinary delta-metrics.  Therefore
\[
          \overline S
          \simeq
          \overline S_m\otimes
          (\overline H^{\otimes m})^\vee
\]
is an ordinary delta-metric by tensor and dual closure.

Now let \(\overline L\) be any cubically metrized line bundle on \(A\).
The line bundle \(L_+\) of \eqref{eq:p09p2-plus-minus} is symmetric, so the
preceding paragraph makes \(\overline L_+\) an ordinary delta-metric.
Equation \eqref{eq:p09p2-delta-equivalence} then makes \(\overline L\) an
ordinary delta-metric.  Thus
\cref{itm:p09p2-ample-symmetric}\(\Rightarrow\)\cref{itm:p09p2-all}, and
the proof is complete.
\end{proof}
